\begin{abstract}
Congressional redistricting requires that districts be drawn with ``substantially equal'' populations, yet the Constitution and Supreme Court precedent leave unresolved the foundational question of which population to equalize. Five distinct definitions compete in law, politics, and administrative practice: total population, voting-age population (VAP), citizen voting-age population (CVAP), registered voters, and actual voters. The Supreme Court's 2016 decision in \emph{Evenwel v. Abbott} declined to mandate any alternative to total population for federal districts, preserving state discretion while foreclosing nothing. We survey the legal landscape from \emph{Reynolds v. Sims} (1964) through \emph{Evenwel} and analyze the practical divergence between these definitions across all 50 states using 2020 Census and American Community Survey data. In 12 states, citizen VAP diverges from total population by more than five percentage points, a difference sufficient to shift district boundaries and partisan outcomes. We introduce the \texttt{--balance-metric} flag in the BISECT redistricting platform, which allows algorithmic redistricting to operate under any of the five definitions, enabling researchers, litigators, and redistricting commissions to test population-definition sensitivity without altering the underlying graph-bisection algorithm. The choice of balance metric is not politically neutral: shifting from total population to citizen VAP systematically advantages rural, predominantly white, low-immigration communities. Transparency about this choice is essential for any redistricting system claiming algorithmic objectivity.
\end{abstract}
